\section{Seed Quality Assessment: Traditional Methods}

Conventional seed quality assessment rests on two broad approaches: manual visual inspection and laboratory-based biochemical testing. In field practice, inspectors evaluate seeds by examining size, shape, colour and surface texture; judgements that are inherently subjective, vary between observers and resist standardisation across districts or seasons. Laboratory tests such as the Tetrazolium (TZ) test and standard germination assays offer more reliable viability estimates, but both are destructive and typically require four to seven days under controlled conditions. That turnaround makes them impractical for screening large consignments at intake points, or for farmers who need a decision before the planting window closes.

Near-infrared (NIR) spectroscopy and X-ray transmission imaging have been investigated as non-destructive alternatives. NIR detects internal moisture gradients and biochemical composition without damaging the sample; X-ray imaging can reveal embryo defects invisible on the seed surface. Both perform well in laboratory conditions, but equipment costs typically USD 15,000-80,000 per unit and the requirement for trained operators confine their use to national seed-testing stations. Smallholder farmers and district-level inspectors in regions such as northeast India have no practical access to either technology, which is precisely where the need for affordable, rapid screening is most acute.

\section{Deep Learning for Agricultural Image Classification}

Convolutional Neural Networks (CNNs) have been applied to a range of agricultural classification problems over the past decade, including leaf disease identification~\cite{mohanty2016plant}, weed recognition~\cite{dyrmann2016}, and post-harvest quality grading~\cite{zhang2020}. Their principal advantage over classical machine-learning pipelines is that spatial features are learned directly from raw pixel data during training, removing the need to manually engineer feature extractors for each new crop or defect type.

Applying CNNs to seed classification follows naturally from this body of work. Paddy seeds exhibit subtle differences in surface texture and colour between certified and counterfeit lots; differences that resist quantification through handcrafted descriptors but emerge reliably from sufficiently large training sets. The computational cost of running a trained CNN at inference time is also modest enough to support server-side web deployment, which matters when target users have intermittent electricity but reasonable mobile-data coverage.

\section{Transfer Learning and ResNet Architectures}

Training a CNN from random initialisation requires on the order of tens of thousands of labelled examples per class to achieve stable convergence. Annotated seed image datasets of that scale do not yet exist for most crop varieties, and certainly not for the specific task of detecting counterfeits. Transfer learning is the standard response: a model pre-trained on ImageNet (1.28 million labelled images across 1{,}000 categories) is fine-tuned on the smaller domain-specific dataset. The early convolutional layers, which encode low-level features such as edges and textures, are retained; only the later task-specific layers are updated. In practice, this approach consistently reaches accuracy close to what a from-scratch model would achieve on a dataset roughly ten times larger.

ResNet architectures are well-suited to this fine-tuning scenario. The residual skip connections introduced by He et al.~\cite{he2016deep} allow gradients to propagate through very deep networks without vanishing, making it feasible to fine-tune models with 50 or more layers on datasets that would cause shallower architectures to overfit. ResNet-50 has become a common baseline in agricultural image classification, balancing parameter count ($\approx$25 million) against accuracy on standard benchmarks, with pre-trained weights publicly available. For this project, it provides a well-validated starting point that reduces both training time and the risk of overfitting on the available paddy seed dataset.

\section{Farmer Feedback Systems}

A detection model trained on a fixed dataset will degrade over time if counterfeiting practices evolve and no new training data are collected. Field-level feedback from farmers is one practical mechanism for capturing these shifts. When a farmer submits a seed batch for image-based verification and subsequently reports germination failure or abnormal yield, that report constitutes a labelled example, a model prediction paired with a ground-truth outcome, that can be incorporated into periodic retraining cycles.

Feedback also serves a more immediate diagnostic function. Aggregating reports across users can flag a suspect lot or supplier before the model itself is updated: three independent failure reports from the same district within a fortnight is a meaningful signal, regardless of what the image classifier predicts. This rule-based alerting complements the model rather than replacing it. Sustaining farmer participation remains a practical difficulty; deployments in low-resource contexts consistently find that single-tap rating interfaces produce far higher response rates than free-text fields or structured forms~\cite{munoz2018ict}.

\section{Web-Based Agricultural Platforms}

Deploying a trained model as a web service separates the computational workload from the end user's device. A farmer uploads a seed image through a browser; the server runs inference and returns a result. This architecture offers two concrete advantages over distributing a standalone mobile application: the model can be updated centrally without requiring client-side installations, and every prediction is logged in a structured format that supports monitoring and retraining.

For farmers in Manipur and comparable regions, browser-based access via a mid-range Android handset is more realistic than assuming a dedicated application installation. Keeping the frontend lightweight, avoiding heavy JavaScript bundles and large media assets reduces page-load times on 4G connections with variable signal quality. Containerised inference endpoints handle computationally intensive tasks without requiring a permanently running server, reducing operating costs during off-season periods when query volume is low.

\section{Review of Related Work}

Early automated seed classification employed Support Vector Machines (SVMs) and Random Forests operating on handcrafted features derived from shape measurements, colour histograms, and texture descriptors. Ni et al.~\cite{ni2009} classified rice seed varieties using morphological parameters extracted from binary silhouettes, achieving 94\% accuracy across five varieties under controlled lighting. These methods perform adequately when intra-class variation is low and imaging conditions are uniform, but generalise poorly to photographs taken in field settings.

Deep learning approaches have largely displaced handcrafted methods in more recent work. Liu et al.~\cite{liu2020} applied VGG16 transfer learning to maize seed defect detection, reporting 97.3\% accuracy on a 12-class dataset. Xu et al.~\cite{xu2021} used ResNet-50 for soybean quality grading and found that fine-tuning from ImageNet weights outperformed training from scratch by approximately 8 percentage points when training data were limited to fewer than 2{,}000 images per class. Both studies, however, address naturally occurring defects, such as cracking, mould and insect damage, rather than intentional substitution. A counterfeit seed is engineered to resemble the genuine article; distinguishing the two is a harder classification problem than separating a damaged kernel from an intact one. The visual gap between classes is narrower, and the adversarial nature of counterfeiting means the distribution of counterfeits may shift as producers adapt. No published study directly evaluates CNN performance on verified counterfeit-versus-genuine seed pairs, and none integrates a structured feedback collection mechanism into the evaluation pipeline.

\section{Summary and Research Gaps}

The reviewed literature establishes that transfer-learned CNNs are effective for agricultural image classification, including seed quality assessment. Three gaps are directly relevant to this work. First, no published study has targeted intentional counterfeiting as distinct from natural quality defects; the two problems differ in that counterfeits may pass casual visual inspection and require the model to detect subtle discriminative cues. Second, existing models are evaluated on laboratory-quality images collected under controlled conditions; performance on photographs taken by farmers using mobile devices in variable lighting has not been reported. Third, no end-to-end system has been published that connects model inference to a structured feedback collection mechanism and uses that feedback iteratively for retraining.

This project addresses all three gaps. The classifier is trained and evaluated specifically on a fake-versus-genuine paddy seed dataset; the web platform is tested with mobile-captured images; and the farmer feedback module is built directly into the application, with logged responses stored in a format suitable for future retraining rounds.